\documentclass[]{article}
\usepackage[]{amsmath}
\usepackage[]{hyperref}
\usepackage[]{smartref}
\usepackage[]{cleveref}
\usepackage[]{marginfix}
\usepackage[]{graphicx}
\bibliographystyle{plain}
\usepackage[a4paper, total={6.5in, 8in}]{geometry}
\usepackage[]{setspace}

\singlespacing

% 
\graphicspath{{/Users/trevorlong/Documents/MIT/PhD/figures/}}

% custom commands
\newcommand{\plant}{$f\left(\mathbf{x},\mathbf{u}\right)$}


\title{Longitudinal Landing Control for Blown Lift Aircraft}
\author{Trevor Long}

\begin{document}
	
	\maketitle
	
	\section{Introduction} \label{sec:introduction}
	
	\begin{figure}[!ht]
		\centering
		\includegraphics[width=0.6\textwidth]{generic-pic.jpeg}
		\caption{(TO DO, PLACEHOLDER PICTURE)Comparison of a common traditional aircraft design (left) and a distributed electric propulsion, blown-lift aircraft (right)} \label{fig:STOL-designs}
	
	\end{figure}
		For the past decade, significant effort has been expended in the aviation industry to design new aircraft that would take advantage of advances in battery storage and electric propulsion technologies to provide new capabilities to open opportunities to serve both dense urban locations in cities as well as less-developed rural areas. 
		At present, these efforts have resulted in significant interest in a new generation of fixed-wing, Short Takeoff and Landing (STOL) aircraft which are able leverage Distributed Electric Propulsion (DEP) systems to take advantage of beneficial aero-propulsive coupling, 'powered-lift', designs \cite{courtin_performance_nodate};
		STOL aircraft provided distinct advantages over Vertical Takeoff and Landing (VTOL) aircraft because while they are able to operate with similar ground infrastructure requirements, they retain most of the cruise efficiency of advantage that conventional fixed-wing aircraft have over VTOL aircraft like helicopters.
		At present several companies such as Electra.aero \cite{electra-first-flight} and Regent \cite{regent-sea-trials} are running flight-tests of full-scale prototypes, and both are actively working on the design and certification of commercial-scale aircraft. \\
		\\
		% Current Status / State of the art
		Electrified STOL aircraft, referred to now comonly as ultra-STOL or uSTOL aircraft, offer significant capability improvements over previous STOL designs due to the use of DEP. 
		\begin{itemize}
			\item The electra EL-2 prototype is able to take-off and land in ~150ft \cite{electra_2023}, more than previous examples \cite{courtin_performance_nodate}
			\item More effectively distribute blowing across the wing, providing more effective blowing \cite{long_experimental_nodate,agrawal_wind_2019,spence_lift_1958}
			\item utilize fly-by-wire flight control systems to augment pilot capabilities \cite{courtin_performance_nodate}
		\end{itemize}
		However, the short-field performance of uSTOL is still largely connected to the management of the stability and control issues associated with low-speed flight \cite{courtin_performance_nodate}.  
		\begin{itemize}
			\item STOL aircraft suffer from poor longitudinal stability and sluggish roll and yaw response to control inputs, these can all be connected to both blowing and slow flight speed and so are endemic to the platform
			\item
		\end{itemize}
		There is also a question of 'how' to fly them because dimension of control is much higher than the number of control axes and the intense coupling between blowing and control authority along the wing.
		\begin{itemize}
			\item Courtin 'fixed' this by tying together inner and outer propellers into units of propellers which were controlled together \cite{courtin_performance_nodate}
			\item with propeller-based blowing there are potential stall regions behind the prop \cite{long_experimental_nodate,courtin_flight_2020}
			\item Power to different motors can be used to control roll/yaw but due to coupling must be done carefully \cite{courtin_performance_nodate} (e.g. might run out of thrust for a certain operation OR too little thrust might cause stall)
		\end{itemize}
		Urban environments are anything but nice to operate in
		\begin{itemize}
			\item urban heat islands contribute to increased thermal activity
			\item large buildings break up steady winds and shed vortices which will make landing difficult
			\item together with the results from courtin it is clear that the robust realization of short landings will require further investigation.
		\end{itemize}
		The reliable operation of uSTOL aircraft is dependent on the development of suitable landing trajectories and control methods which decrease the maximum landing dispersion. 
		\begin{itemize}
			\item requires we know how different disturbances effect the aircraft
			\item requires some knowledg of the 'best' way to respond to 1) design controllers and 2) grade their performance
			\item will most likely require the development of pilot aids either through autoland systems or control configurations which offload pilot
		\end{itemize}
		Advances in control techniques from the robotics field provide interesting avenues of investigation \cite{tedrake_underactuated_nodate} \cite{nguyen_model_2020}.
		\begin{itemize}
			\item Moore was able to perch a glider using simple model and limited control using trajectory optimization to design trajectory and LQR control to manage deviations from that trajectory \cite{moore_robust_nodate}
			\item More recent investigations have used model predictive control \cite{moore_precision_2023}
			\item can leverage these techniques to help control this problem which is much more complex than the standard aircraft control problem
		\end{itemize}
		
		\subsection{Thesis Outline and Objectives}
		%Thesis goal
		The objective of this thesis is to develop methods to limit the landing dispersion of uSTOL aircraft. Specifically, this thesis will focus on control of the longitudinal problem adapted from \cite{courtin_performance_nodate} with the addition of disturbances during landing approach shown in \cref{fig:landing-problem}. 
		The investigation approach intends to follow methods used by \cite{moore_precision_2023,moore_robust_nodate} to develop state and control trajectories with robust disturbance rejection characteristics. 
		It will then explore possible control techniques, and sensor suites required for this control. 
		The powered-lift plant for all of the above will come from a combination of low-fidelity, mid-fidelity, and wind-tunnel based aerodynamics databases. 
		Specific thesis objectives are outlined below: 
			  
		\begin{figure}[!t]
			\centering
			\includegraphics[width=0.6\textheight]{landing-problem.pdf}
			\caption{The premise of the landing problem} \label{fig:landing-problem}
		\end{figure}
		\begin{itemize}
			\item Investigate the effects of gust disturbances on the stability of blown-lift SSTOL aircraft and identify effects which lead to trajectory instability
			\item Development of a simulation environment whcih can be used to investigate SSTOL control problems through a trajectory optimizaiton framework
			\item Development of detailed aerodynamics databases at mutliple fidelities for investigation
			\item Development of landing trajectories that have robust properties and control authority for gust rejection. 
			\item investigation of methods which can be used to aid pilots in flying these trajectories 
		\end{itemize}

	In the remainder of this document more context on issues and methods referred to in this section.
	\newpage 
	
	\section{Background} \label{sec:problem-statement}
	% briefly introduce STOL
	
	Despite the current generation of blowing-augmented STOL aircraft being designed to take advantage modern electric propulsion and battery technologies, many attempts at successful designs using this concept have been proposed since the jet-flap concept was proposed by Spence \cite{spence_lift_1958}. 
	The jet flap allows for significant increases in ahievable lift coefficients, far beyond what is known to be achievable using multi-element flaps alone \cite{smith-high-lift,long_experimental_nodate,long_parametric_2021}
	Previously these designs were explored extenstively in the 1970s and 1980s, several groups were experimenting with different variations of the concept; however, none of these designs ultimately progressed past the test phase even with successful test campaigns; the common issue cited for the ending of these programs was the mechanical complexity required to control and design these aircraft. Courtin provides a detailed overview of the history of these aircraft is provided by Chapter 2 in \cite{courtin_performance_nodate}.
	\\
	The obvious advantage of generating high lift coefficients is the ability to reduce the footprint of landing and takeoff. While a modern aircraft using flaps and slats might generate a $c_\ell \approx 2$, a well-designed blowing-augmented aircraft might be able to easily reach a $c_\ell \approx 5$ \cite{smith-high-lift,long_experimental_nodate}; All else being equal, the STOL aircraft using blowing would then fly $\approx 1.6x$ slower than the conventional aircraft. Taken to the extreme it seems likely that for a Cesna-sized aircraft it is possible for a fixed-wing aircraft to take-off and land along a $150$ft runway which would allow operations in both highly urban and remote environments \cite{courtin_performance_nodate}.\\
	\\
	But as it turns out, flyig slowly and landing repeatably is more of a challange than it first may appear. 
	Additionally, the aerodynamics and rigid dynamics of these aircraft are much more complicated than those for conventional aircraft designs. 
	Additionally, the use of distributed electric propulsion and fly-by-wire systems gives designers now hav emore nuanced control on the wing. 
	In the remaineder of this section, I will briefly summarize many of the issues present in 
	
	% what makes sstol aircraft different from a 'normal' aircraft
	\subsection{Peculiarities of Distributed Electric STOL Aircraft}
Objective: What about the aerodynamics and dynamic system of a STOL aircraft makes it more complicated than a normal plane
\subsubsection{Aerodynamics} \label{sec:aero}
- what is STALL? does it exist?
- control reversal / steep performance curves
- extremely coupled with thrust, lift, pitching moment
- extreme circulation
- this has ONLY to do with the 2D 3D aerodynamics, not anything about control, dynamic pressure, or any of those things

The aerodynamics and performance of blowing-augmented airfoils is much different than those for clean wings. 
The starting point for this is the understanding that while the performance of a 'clean' 2-dimensional wing with a single flap, which will define as $c_\ell, c_m, c_d$ are $f(\alpha, \delta_f, Re)$ while for a blown wing we now must consideer a new blowing parameter giving the performance as $f(\alpha,\delta_f, Tc',Re)$ \cite{spence_lift_1958}. 
Thus aerodynamics and thrust have become coupled and as \cite{courtin_performance_nodate} observes, this can make it difficult to find an operating point with net-0 x-force for landing. \\
\\
The airfoil dynamics themselves are also complex. 
Historically we prefer to fly aircraft in the linear regions of the performance curves where performance is predictable within an acceptable margin \cite{drela_flight_2014,etkin_dynamics_1959}, but for blowing-augmented aircraft this likely is not possible if we want to use them to their fullest potential. 
The performacne of these airfoils is very non-linear especially in high-lift regions as seen in \cref{fig:blown-airfoil-performance} and it seems unlikely that traditional flight controllers are suited to handling these problems.
sign reversals and large slopes can be seen in $c_m$, and even the definition of stall is hard to define exactly for these airfoils \cite{long_parametric_2021}.\\
\\

\begin{figure}{!h}
	\centering
	\includegraphics[width=0.8\textwidth]{blowing-performance.png}
	\caption{Performance curves showing the primary $c_\ell, c_x, c_m$ relationships for a blown wing as a function of $\alpha$, $\delta_f$, and $\Delta c_J$ from \cite{long_parametric_2021}. Aerodynamic performance is shown  to be extremely nonlinear along different control dimensions} \label{fig:blown-airfoil-performance}
\end{figure}

\subsubsection{Low Flight Speeds and Disturbance Sensitivity}
- one is the fact that we fly really slow
- look at courtin thesis for exact sensitivites, he did a section on this
- know from experience roll and pitch are super sensitive
- this is the motivation to look at gusts

Slow flight speeds and the nonlinear dynamics make these aircraft much more vulnerable to disturbances. While a typical Cessna 172 might land at around $60$kts tof airspeed, a blowing-augmented on might be able to do so at nearly half that $30kts$ \cite{courtin_performance_nodate}. 
Lower flight speeds make the aircraft much more sensitive to gusts in the nonlinear manner as shown in \cref{eq:gust_effect} for an aircraft in level flight encoutnering an updraft.
Thermal drafts and shears are very common near airports and more common in cities (ref needed), additionally in areas with high buildings or peaks and inconsistent ground heating much more unpredictable and much stronger than those in more consistent conditions. 
On approach the potential for large swings in $\alpha$ due to gusts, shears, and other disturbances combined with the nonlinear aerodynamics from the previous section indicate that these aircraft will be difficult to control in 'real' environments. 
\begin{equation} 
	\Delta\alpha_\mathrm{gust} = \arcsin\left(\frac{w_g}{V_\infty}\right) \label{eq:gust_effect}
\end{equation}
\begin{figure}[!h]
	\centering
	\includegraphics[width=0.8\textwidth]{generic-pic.jpeg}
	\caption{Change in $\alpha$ as a function of true airspeed. Operating points for SSTOL aircraft feel much larger effect for the same-sized gust} \label{fig:gust-effect}
\end{figure}

\subsubsection{Control Saturation} \label{sec: saturation}

%	- flying at low speeds is great, but there is tension between flying slow and having enough control authority to do anything
%	- at low speeds any unblown surface is going to have less authority
%	- additionally, with strong pitching moment, it is possible
For these aircraft designs there is a tension between the desire to blow more and fly more slowly and aircraft controllability.
The issue is that as more blowing is added to the system and the aircraft is able to achieve higher lift coefficients and thus fly more slowly, unblown surfaces on the aircraft experience both less dynamic pressure and significant influence from the circulation generated by the blown wing. 
Thus even with reasonably-sized control surfaces and throws on controls such as the elevator, it is extremely likely that errant disturbances or control inputs could cause a problematic loss of control. \\
\\
The most apparent reaosn for control saturation is the loss of dynamic pressure on unblown control surfaces. 
At high blowing setttings the velocity ratio $\frac{V_j}{V_\infty}$ have been observed to easily be as high as 3 \cite{long_parametric_2021} and will most likely be even higher for production aircraft \cite{courtin_performance_nodate}. 
It is readily apparent from the definition of dynamic pressure that the ratio between blown and unblown surfaces will obey a square law, so that the observed ratio $\frac{V_j}{V_\infty} = 3$ will lead to a dynamic pressure ratio of $\frac{q_\mathrm{blown}}{q_\mathrm{unblown}} = 9$. 
This leads to two connected issues: for short landing we want to push these aircraft to fly as slow as possible, but at some point there crossees a line where unblown controls simply do not have enough dynamic pressure and are likely to either simply not generate enough force to balance those from the wing or stall while doing so leading to the loss of control. \\
\\
This perhaps would not be a significant issue if surfaces were not signficiantly influenced by the wing, but for blown aircraft the exact opposite is true. 
While it is well known that even for conventional aircraft the wing imparts an $\Delta\alpha$ at the tail this is for the most part a very miinor impact for lift coefficients we tend to fly at (in the neighborhood of 1-2). 
Blown lift aircraft fly at much higher lift coefficients and thus from theory alone we would expect similarly high $\Delta\alpha$ around the wing from induced velocities: multiple sources confirm this to be the case \cite{courtin_performance_nodate} \cite{courtin_flight_2020} (need a few more sources here).
Large $\alpha$ offsets for blowing operating points imply that a significant control effort must be applied to simply stabilize the aircraft, most acutely along the pitch axis; such large control efforts may push these surfaces closer to stall or saturation. 
Thus it is possible to enter a situation where a small disturbance in the biased control direction could push a surface past its authority limits by either stalling it or hitting a throw limit leading to loss of control. \\
\\
Lastly, another peculiarity of blown lift aircraft is the large slopes observed in pitching moment and the sign reversal on coefficients at certain operating points. 
This is particularly apparent in the $c_m-\alpha$ curves for a 2-dimensional wing shown in \cref{fig:blown-airfoil-performance}, where for all blowing cases a steep slope and sign reversal is observed near $\alpha = 15^\circ$ for all flap deflections \cite{long_parametric_2021}. 
These control reversal can be extremely hazardous as it is possible that small perturbations can push the aircraft operating points suddenly and without warning. 
Limitations on actuation speed, as well as the bias and saturation limits described earlier in this section, mean that if these disturbances were to happen quickly enough or are large enough, it is possibel to enter an uncontrollable operating point. 
While we might seek to simply stay away from operating regions with large slopes and sign reversal, these regions also happen to overlap with high-$c_\ell$ operating regions which are desireable for STOL performance.\\
\\

\subsubsection{High-Dimensional Control Space} \label{sec:acuation}
- what actuator do you use and do respond to what?
- there are many more combiantions than previously
- control is incredibly coupled
- hard limits on some, do you use others to compensate?
For the direct control of blown-lift aircraft we can make two observations based on \cref{fig:courtin-control-layout}. The first is that the number of possible control inputs is much higher than on a conventional aircraft due to the number of propellers and flaps which can be actively managed; and the second is that each of these has a much more complex dynamic response than conventional counterparts. 
For example, a change in flap setting results in a constant shift in $C_L, C_M, C_D$ and thurst can be modulated to maintain speed. 
But based on the discussion in \cref{sec:aero}, it is quickly seen that thrust modulation will now also impact $C_L, C_D, C_M$.
The influence of this change is also extremely location dependent; flaps and propellers on outer edges of the wing have more roll impact than those more center, and the effect of induced velocity is also dramatic \cite{courtin_performance_nodate}.\\
\\
There is also the question of how-best to use the large number of contorls. 
A conventional aircraft has flight controls that each have a very clear primary axis of effect: ailerons control roll, rudder controls yaw, elevator controls pitch, throttle controls velocity, and so on. 
Many acutators gives us many possibilities and while the system is still technically under-actuated based on the definition of Chapter 1 of \cite{tedrake-underactuated}, it is very non-unique.
This suggests the question of using controls and control combinations as effectively as possible. \\

Courtin suggests control groupings as a means of handling redundency in the control space \cite{courtin_performance_nodate}.
However, it is very likely that for many maneuvers there are more effective or safe combinations of control inputs that result in the same dynamic response and with modern fly-by-wire systems control gains are easy to change.
How best explore the space of possible control inputs for different maneuvers is a question that this thesis will explore.\\
%	\cite{under-actuated robotics here}
%	over-actuated because we have many of the same control
%	- normal aircraft is nicely split so we typically have 1, maybe 2 controls per-axis for redundency
%	- think 10 propulsors so we are over-actuated on thrust
%	- but we can also see that the same applies for roll
%	-courtin tried to solve this by reducing control dimensionality and grouping controls, but this leaves out opportunity to do better
%	- non-uniqueness is the central theme, and we want to sift through the trajectory design space for 'best' solutions
\begin{figure}[!h]
	\centering
	\includegraphics[width=0.5\textwidth]{electra-el-2.png}
	\caption{Control layout for a distributed electric propulsion aircraft from Courtin} \label{fig:courtin-control-layout}
\end{figure}
	% why is this an issue?
	\subsection{Wind and Landing} \label{sec:winds}
Objective: overview for why winds during landing are bad and why cities make the problem worse


- what are the disturbances that they will encounter flying?
- what does an airport currently look like?
- how is this any different?
- what level of disturbance do we care about?
- honeslty the detail here isn't important, we are simply showing that even small disturbances are incredibly meaningful here

Low-and surface level winds have a significant effect on landing for aircraft. 
\begin{itemize}
	\item bulk air movements (steady winds)
	\item turbulence (from storms, unstable air)
	\item local structures and thermals
\end{itemize}
While they affect both runway operations and landing procedures, we'll focus on the potential impact on aircraft. \\
\\
In particular we are worried about two things. Turbulence generated by ground structures (why airports tend to have low buildings for flight restriction) and thermal activity on the ground. 

Turbulence generation in urban environments

Thermal updrafts, effect of urban environment on updrafts (hotter -> more energy)

Why it matters more for DEP STOL aircraft -- bring it back to the AoA / dcj diagram an nonlinear aero
	% How do you figure out how to control them 
	\subsection{Trajectory Planning} \label{sec:trajectory-planning}
How are trajectories planned now (hint, they've been developed over time by the FAA). methods are inadequate for much higher performance aircraft. 

one way to look at this is as a trajectory optimization problem, now talk some about that


\begin{equation} \label{eq:traj-optimization-discrete-time}
	\begin{array}{rrclcl}
		\displaystyle \min_{\mathbf{x}[k],\mathbf{u}[k]} & \multicolumn{3}{l}{ \ell\left( \mathbf{x}[N],\mathbf{u}[N]\right) + \sum_{k=0}^{k=N} J\left( \mathbf{x}[k],\mathbf{u}[k]\right)}\\
		\textrm{s.t.} & \mathbf{x}[k+1] = f\left( \mathbf{x}[k], \mathbf{u}[k]\right)  \\
		& g(\mathbf{x}[k] \mathbf{u}[k])\leq 0  \\
		& h(\mathbf{x}[k],\mathbf{u}[k] ) = 0\\
		& + \; \mathrm{other \; constraints}...\\
	\end{array}
\end{equation}

What is optimal control / trajectory optimization? What was it developed for and why hasn't it been used for aircraft? 
- historically you could easily separate controls to individual components
- they've used it, but not to design control trajectories (GPOPS, courtin, soaring)
- how is the problem defined (tedrake, other AIAA paper that described collocation)
\\
How exactly do you set this up as a trajectory problem? What is the problem, what are the constraints?\\
\\
ARe there any examples of similar funcitoning systems? (perching) \\
\\
	% Once you know how do you control / aid landing
	\subsection{Aircraft Control Systems} \label{sec:control}
\begin{itemize}
	\item Modern fly-by-wire systems, make control mixing easier can switch between 'modes' more easily
	\item All sorts of aids: autopilot, stability enhancement, envelope control. Ultimate goal is to have one or multiple of these systems
	\item The plan is to understand how to fly with the trajectory optimization, then implement as either a controller or flight aid, whichever seems more relevant
	\item generate procedure -> generate controller
\end{itemize}

Since at least the 60s we have also had autoland systems

\subsubsection{Trajectory Planning} \label{ssec:trajectory-planning}
First thing we can do is use modern tools to generate trajectories for these aircraft using surrogate models. By doing this we can use the optimizer to help us understand the difference between 'good' control and 'bad' control. This is especially relevant because of how high the control space is which makes it much more difficult to decide what conrols what. \\
\\
Why is this relevant now and not before?\\
\\
Why do we have confidence we can come to some actual solution\\
\\
To do this, what is the plan?\\
\\
How complete does the model need to be?\\
\\
\begin{itemize}
	\item technique of trajectory optimization is nothign novel, GPOPS, others have been doing it for a while
	\item lots of advancement for using on-line in robotics, particularly for manipulation
	\item example: used for perching paper as exmaple of solving complex aerodynamic model
	\item h
\end{itemize}

what is novel here is proposing to develop control trajectories, not just aircraft tajectories. 

\subsubsection{Implementation} \label{ssec:controller}
Once you generate flight paths you have the option of how you use that information. \\
\\
Do you create a controller and try and eliminate pilot ?\\
\\
Do you create pilot aids like alarms and vocal warnings? maybe something more modern \\
\\
Do you try and use the fly-by-wire system to create a dew control schema that the pilot switches to?\\
\\
How does this impact operations more broadly?\\
\\



	
	\bibliography{bibliography}
	
	
\end{document}